%% ============================================================================
%% Manuscript prepared for: Biomedical Signal Processing and Control (Elsevier)
%% Template: Standard Elsevier elsarticle (falls back cleanly to article if elsarticle.cls is not installed locally)
%% ============================================================================
\documentclass[11pt,a4paper]{article}
\usepackage[utf8]{inputenc}
\usepackage[T1]{fontenc}
\usepackage[english]{babel}
\usepackage[margin=2.3cm]{geometry}
\usepackage{amsmath,amssymb,amsfonts}
\usepackage{booktabs}
\usepackage{graphicx}
\usepackage{hyperref}
\usepackage{microtype}
\usepackage{xcolor}
\usepackage{tabularx}
\usepackage{multirow}
\usepackage{array}
\usepackage{float}
\usepackage{caption}

\hypersetup{
    colorlinks=true,
    linkcolor=blue!60!black,
    citecolor=blue!60!black,
    urlcolor=blue!60!black
}

\title{\LARGE \textbf{Hybrid Multimodal Deep Learning and Handcrafted Dermoscopic Feature Fusion with Lesion-Grouped Validation and Clinical Threshold Optimization for Multiclass Skin Cancer Classification}}

\author{
    \textbf{Alexandre Ramos da Silva}\\[0.3em]
    \small Postgraduate Program in Data Science and Artificial Intelligence in Healthcare, Brazil\\
    \small \texttt{github.com/aramos197442/skin\_cancer}
}
\date{\today}

\begin{document}
\maketitle

%% ============================================================================
%% ELSEVIER HIGHLIGHTS (Strictly <= 85 characters per bullet including spaces)
%% ============================================================================
\noindent\rule{\textwidth}{0.6pt}
\subsection*{Highlights}
\begin{itemize}
    \item Lesion-grouped 70/15/15 split eliminates multi-image data leakage in HAM10000.
    \item Systematic ablation from handcrafted ABCD descriptors to hybrid deep learning.
    \item Hybrid EfficientNetV2B0 + ABCD fusion achieves 76.6\% accuracy and 0.603 MCC.
    \item Post-hoc class threshold tuning improves macro sensitivity (+1.6 percentage points).
    \item Grouped bootstrap, ECE/Brier calibration, Grad-CAM, and perturbation audits.
\end{itemize}
\noindent\rule{\textwidth}{0.6pt}
\vspace{0.5em}

%% ============================================================================
%% STRUCTURED ABSTRACT (Elsevier BSPC Format)
%% ============================================================================
\begin{abstract}
\noindent\textbf{Background and Objective:} Computer-aided diagnosis (CAD) of dermatoscopic images requires both high sensitivity for life-threatening malignancies and rigorous methodological isolation to avoid optimistic bias from duplicate lesion views. This study investigates the progression from interpretable handcrafted dermatoscopic descriptors to a hybrid multimodal deep learning architecture on the seven-class HAM10000 benchmark, explicitly addressing lesion-level data leakage, severe class imbalance, probabilistic calibration, and clinically oriented decision threshold optimization.\\
\noindent\textbf{Methods:} All 10,015 dermatoscopic images (7,470 unique lesions) were partitioned into 70\% training, 15\% validation, and 15\% hold-out test sets using stratified grouping by \texttt{lesion\_id}, ensuring zero lesion overlap across splits. Hair artifacts were mitigated via morphological Black-Hat filtering and Telea inpainting (DullRazor) on $224 \times 224$ center-cropped images. A 30-dimensional clinical feature vector encoding the dermatological ABCD rule---RGB and HSV chromatic means (6), Haralick Gray-Level Co-occurrence Matrix (GLCM, 4) and uniform Local Binary Patterns (LBP, 13), and Otsu-segmented log-transformed Hu invariant moments (7)---was first evaluated incrementally with a class-balanced Random Forest baseline. Subsequently, a hybrid architecture fusing an ImageNet-pretrained \textbf{EfficientNetV2B0} visual backbone with the standardized ABCD tabular branch was trained in TensorFlow/Keras under a two-stage fine-tuning protocol using Sparse Focal Loss ($\gamma=2.0$) and uniform class sampling, and benchmarked against \textbf{ResNet50} and \textbf{ConvNeXt-Tiny} across multiple random seeds. Post-hoc per-class decision thresholds ($\theta_k$) were calibrated exclusively on the validation split.\\
\noindent\textbf{Results:} In the Random Forest baseline, chromatic descriptors alone achieved 44.0\% multiclass accuracy; incorporating GLCM and LBP texture descriptors increased accuracy to 54.9\% (+10.9 percentage points), whereas full ABCD + DullRazor reached 53.3\% multiclass accuracy (70.5\% One-vs-Rest balanced accuracy), constrained by severe manifold overlap observed in t-SNE projections. Transitioning to the threshold-optimized hybrid \textbf{EfficientNetV2B0} elevated hold-out test performance to \textbf{76.6\% Multiclass Accuracy} (95\% grouped bootstrap CI: 74.4\%--78.7\%), \textbf{93.3\% Mean OvR Accuracy}, \textbf{74.6\% Multiclass Balanced Accuracy (Macro Recall)} (95\% CI: 71.2\%--77.8\%), \textbf{0.646 Macro F1-Score}, \textbf{0.603 Macro OvR Matthews Correlation Coefficient (MCC)}, \textbf{0.949 Macro ROC-AUC}, and \textbf{0.670 Macro PR-AUC}. Probabilistic calibration yielded a multiclass Brier Score of 0.284 and an Expected Calibration Error (ECE) of 0.068 (6.8\%).\\
\noindent\textbf{Conclusions:} Within the evaluated lesion-grouped protocol, hybrid transfer learning with EfficientNetV2B0 significantly outperforms classical feature-engineered ensembles ($p < 0.0001$, McNemar and paired cluster bootstrap tests) while maintaining low computational overhead ($5.9 \times 10^6$ backbone parameters). Clinical error auditing of melanoma false negatives, Grad-CAM visual explanations, and photometric perturbation tests underscore both the screening potential of threshold calibration and the need for external multi-center validation prior to clinical deployment.

\vspace{0.6em}
\noindent\textbf{Keywords:} Skin Cancer Classification; HAM10000; Dermoscopy; EfficientNetV2B0; Multimodal Deep Learning; ABCD Rule; Threshold Calibration; Data Leakage Prevention.
\end{abstract}

%% ============================================================================
%% 1. INTRODUCTION
%% ============================================================================
\section{Introduction}
Cutaneous malignancies represent the most prevalent class of human cancers worldwide \cite{sung2021, siegel2024}. Although non-melanoma skin cancers such as basal cell carcinoma (BCC) and squamous cell precursors like actinic keratoses (AKIEC) account for the majority of clinical incidences, cutaneous melanoma (MEL) remains responsible for the vast majority of skin cancer mortality due to its rapid metastatic potential \cite{rigel2010}. Early detection via dermatoscopic inspection raises five-year melanoma survival above 95\%, whereas late-stage diagnosis carries a dismal prognosis \cite{ Kittler2002, tschandl2018}. However, visual inspection of dermatoscopic images is subject to substantial inter-observer variability among clinicians, motivating the development of reliable Computer-Aided Diagnosis (CAD) systems \cite{esteva2017, haenssle2018}.

Automatic multiclass classification on dermatoscopic repositories such as HAM10000 \cite{tschandl2018} poses four intertwined biomedical signal processing challenges:
\begin{enumerate}
    \item \textbf{Acquisition Artifacts:} Cutaneous hair shafts, ruler markings, and air bubbles occlude pigment networks and distort lesion boundaries \cite{lee1997, celebi2019}.
    \item \textbf{Extreme Class Imbalance:} Benign melanocytic nevi (NV) represent approximately 67\% of HAM10000, outnumbering rare categories such as dermatofibroma (DF, 1.1\%) and vascular lesions (VASC, 1.4\%) by more than 45:1 \cite{tschandl2018}.
    \item \textbf{Asymmetric Clinical Misclassification Costs:} A false negative (FN) in melanoma delays potentially curative excision, whereas a false positive (FP) leads to a confirmatory biopsy---a clinically acceptable trade-off in screening triage \cite{codella2019}.
    \item \textbf{Lesion-Level Data Leakage in Literature Benchmarks:} HAM10000 contains 10,015 images acquired from 7,470 unique lesions (\texttt{lesion\_id}), meaning multiple images correspond to sequential or magnified captures of the exact same lesion. Studies that perform naive image-level random splitting inadvertently place views of the same lesion in both training and test sets, artificially inflating reported accuracies above 90\% \cite{cassidy2022, maron2019}.
\end{enumerate}

Historically, dermatoscopic CAD relied on handcrafted descriptors inspired by clinical semiotics, notably the ABCD rule (\textit{Asymmetry, Border, Color, Diameter}) \cite{nachbar1994}, combined with texture operators such as Gray-Level Co-occurrence Matrices (GLCM) \cite{haralick1973}, Local Binary Patterns (LBP) \cite{ojala2002}, and morphological shape descriptors like Hu invariant moments \cite{hu1962}, classified via Random Forests \cite{breiman2001} or Support Vector Machines. Conversely, modern deep convolutional neural networks (CNNs) such as ResNet \cite{he2016}, ConvNeXt \cite{liu2022}, and EfficientNetV2 \cite{tan2021}, as well as Vision Transformers (ViT) \cite{dosovitskiy2021, liu2021swin}, learn hierarchical representations directly from pixels via transfer learning from ImageNet \cite{deng2009}.

Rather than treating handcrafted engineering and deep learning as disconnected paradigms, this work presents a systematic, leak-free experimental investigation comparing progressive feature engineering, component ablation, and hybrid multimodal deep learning under strictly identical lesion-grouped partitions. The main scientific contributions of this study are fourfold:
\begin{enumerate}
    \item \textbf{Strict Lesion-Grouped Protocol and Metric Disambiguation:} We enforce a 70\%/15\%/15\% train/validation/test partition grouped by \texttt{lesion\_id} with zero overlap, and mathematically disambiguate global multiclass accuracy from mean One-vs-Rest (OvR) binary accuracy, preventing optimistic reporting distortions under severe class imbalance.
    \item \textbf{Incremental Feature Engineering and Component Ablation:} We quantify the marginal contribution and collinearity trade-offs of chromatic (RGB/HSV), textural (GLCM/LBP), morphological (Otsu + Hu moments), and hair-inpainting (DullRazor) stages in a Random Forest baseline.
    \item \textbf{Hybrid Multimodal Transfer Learning and Comparative Benchmarking:} We fuse an ImageNet-pretrained \textbf{EfficientNetV2B0} backbone with the 30-dimensional ABCD clinical feature vector under Sparse Focal Loss and balanced mini-batch sampling, benchmarking it against visual-only EfficientNetV2B0, \textbf{ResNet50}, and \textbf{ConvNeXt-Tiny} across multiple random seeds with paired statistical tests and computational profiling.
    \item \textbf{Clinically Oriented Threshold Optimization, Calibration, and Error Auditing:} We formally separate probabilistic calibration (Brier Score, ECE, Temperature Scaling) from post-hoc per-class threshold tuning ($\theta_k$), and conduct an in-depth clinical audit of melanoma false negatives, Grad-CAM explainability, and robustness to photometric perturbations.
\end{enumerate}

%% ============================================================================
%% 2. RELATED WORK
%% ============================================================================
\section{Related Work}
Early dermatoscopic CAD systems translated clinical heuristics---such as Stolz's ABCD rule \cite{nachbar1994} and the 7-point checklist \cite{argenziano1998}---into quantitative pipelines comprising preprocessing (e.g., DullRazor hair removal \cite{lee1997}), lesion segmentation via Otsu thresholding \cite{otsu1979} or active contours, and extraction of colorimetric, GLCM \cite{haralick1973}, and LBP \cite{ojala2002} descriptors \cite{celebi2007, barata2014}. While interpretable, handcrafted features struggle with high intra-class variability across amelanotic melanomas, pigmented basal cell carcinomas, and seborrheic keratoses \cite{barata2018}.

Following the landmark demonstration by Esteva et al. \cite{esteva2017} and the release of the HAM10000 dataset by Tschandl et al. \cite{tschandl2018} within the ISIC challenges \cite{codella2019, tschandl2020}, deep transfer learning became the dominant paradigm. Numerous studies explored ResNet50 \cite{he2016}, DenseNet201 \cite{huang2017}, MobileNetV2 \cite{sandler2018}, EfficientNet \cite{tan2019, tan2021}, ConvNeXt \cite{liu2022}, and Vision Transformers \cite{dosovitskiy2021, aladhadh2022, xin2022} for seven-class HAM10000 classification. Multimodal fusion combining CNN visual embeddings with patient metadata or handcrafted radiomic descriptors has also shown promise in improving minority-class discrimination \cite{yap2018, pacheco2020, ningrum2021}.

Nevertheless, critical methodological audits by Cassidy et al. \cite{cassidy2022} and Bissoto et al. \cite{bissoto2020} revealed that a substantial fraction of published HAM10000 studies suffer from three recurring flaws: (i) performing random image-level splits or applying synthetic oversampling (e.g., SMOTE) prior to partitioning, leaking identical \texttt{lesion\_id} views into the test set; (ii) reporting mean One-vs-Rest accuracy ($\sim 90\text{--}95\%$) as if it were overall multiclass accuracy; and (iii) omitting confidence intervals, probabilistic calibration metrics (ECE/Brier) \cite{guo2017}, and paired statistical hypothesis testing \cite{demvsar2006}. Our methodology is specifically designed to close these gaps.

%% ============================================================================
%% 3. MATERIALS AND METHODS
%% ============================================================================
\section{Materials and Methods}

\subsection{Dataset and Leak-Free Lesion-Grouped Partitioning}
We utilized the complete HAM10000 (\textit{Human Against Machine with 10000 training images}) benchmark \cite{tschandl2018} (ingested via the standardized repository \texttt{marmal88/skin\_cancer}), comprising $N = 10,015$ dermatoscopic RGB images acquired from 7,470 unique lesions (\texttt{lesion\_id}) across seven diagnostic categories: \textit{Actinic Keratoses and Intraepithelial Carcinoma} (AKIEC, 327 images), \textit{Basal Cell Carcinoma} (BCC, 514), \textit{Benign Keratosis-like Lesions} (BKL, 1,099), \textit{Dermatofibroma} (DF, 115), \textit{Melanocytic Nevi} (NV, 6,705), \textit{Melanoma} (MEL, 1,113), and \textit{Vascular Lesions} (VASC, 142). Ground-truth labels in HAM10000 were established via histopathology ($>50\%$), digital follow-up, in-vivo confocal microscopy, or expert consensus \cite{tschandl2018}.

To prevent lesion-level data leakage, the dataset was partitioned into \textbf{70\% Training}, \textbf{15\% Validation}, and \textbf{15\% Hold-Out Test} ($N_{\text{test}} = 1,437$ images) using \texttt{StratifiedGroupKFold} grouped strictly by \texttt{lesion\_id} and stratified by diagnostic label (master seed fixed at 42; multi-seed sensitivity evaluated across seeds $\{42, 100, 2024\}$). Programmatic assertions verified zero intersection of \texttt{lesion\_id} sets across the three splits:
\begin{equation}
    \mathcal{L}_{\text{train}} \cap \mathcal{L}_{\text{val}} = \varnothing, \quad
    \mathcal{L}_{\text{train}} \cap \mathcal{L}_{\text{test}} = \varnothing, \quad
    \mathcal{L}_{\text{val}} \cap \mathcal{L}_{\text{test}} = \varnothing
\end{equation}
All feature selection (\texttt{VarianceThreshold}), feature standardization (\texttt{StandardScaler}), class weighting, and synthetic resampling were fitted \textbf{strictly on the training split}. Model selection (early stopping, learning rate reduction) and post-hoc threshold calibration ($\theta_k$) were performed \textbf{exclusively on the validation split}, keeping the hold-out test set completely isolated until final evaluation.

\subsection{Image Preprocessing and Hair Removal (DullRazor)}
Each original dermatoscopic image ($600 \times 450$ pixels) first underwent a 1:1 center crop ($\min(W, H)$) to preserve the native aspect ratio and circular morphology of the skin lesion, followed by bilinear resizing to $224 \times 224$ pixels. To attenuate dark hair shafts that corrupt texture and boundary descriptors, we implemented the \textbf{DullRazor} algorithm \cite{lee1997}:
\begin{enumerate}
    \item Conversion of the RGB image to 8-bit grayscale ($I_{\text{gray}}$).
    \item Morphological \textbf{Black-Hat} filtering using a $9 \times 9$ rectangular structuring element $B_{9\times 9}$: $I_{\text{bh}} = (I_{\text{gray}} \bullet B_{9\times 9}) - I_{\text{gray}}$, isolating thin dark tubular structures.
    \item Binary thresholding ($I_{\text{bh}} > 10$) to obtain the hair occlusion mask $M_{\text{hair}}$.
    \item Fast marching restoration of occluded pixels via Telea's inpainting algorithm (\texttt{cv2.INPAINT\_TELEA}, radius $= 1$ px) \cite{telea2004}.
\end{enumerate}

\subsection{Handcrafted Clinical Feature Extraction (30 ABCD Descriptors)}
To translate dermatological semiotics into a compact tabular representation, we extracted a 30-dimensional feature vector $\mathbf{f} \in \mathbb{R}^{30}$ structured into three clinical modules (matching the exact implementation in \texttt{eda\_pca\_skincancer.ipynb}):
\begin{enumerate}
    \item \textbf{Chromatic Module (C --- 6 descriptors):} Mean pixel intensities across the three normalized RGB channels ($R, G, B \in [0, 1]$) and the three HSV channels ($H, S, V$), decoupling chromaticity from dermatoscope illumination intensity.
    \item \textbf{Border and Texture Module (B --- 17 descriptors):}
    \begin{itemize}
        \item \textit{Haralick GLCM (4 descriptors):} From the 256-level grayscale co-occurrence matrix computed at distance $d=1$ pixel and angle $\theta=0^\circ$ (symmetric, normalized) \cite{haralick1973}, we extracted \textit{Contrast}, \textit{Homogeneity}, \textit{Energy}, and \textit{Correlation}.
        \item \textit{Uniform Local Binary Patterns (13 descriptors):} Using the circular operator $\text{LBP}_{P=8, R=1}^{\text{uniform}}$ \cite{ojala2002}, we computed a normalized 13-bin density histogram capturing micro-textural patterns (spots, flat areas, edges, and corners).
    \end{itemize}
    \item \textbf{Asymmetry and Morphology Module (A --- 7 descriptors):} The lesion core was automatically segmented from the grayscale image via Otsu's adaptive thresholding \cite{otsu1979} (with polarity correction ensuring background $= 0$). From the resulting binary mask, we computed the seven scale-, translation-, and rotation-invariant \textbf{Hu Moments} ($h_1, \dots, h_7$) \cite{hu1962}, stabilized via a signed logarithmic transform:
    \begin{equation}
        \tilde{h}_i = -\operatorname{sgn}(h_i) \log_{10}\!\left(|h_i| + 10^{-12}\right), \quad i \in \{1, \dots, 7\}
    \end{equation}
\end{enumerate}
Zero-variance checks (\texttt{VarianceThreshold}) and $Z$-score standardization (\texttt{StandardScaler}) were fitted on $\mathbf{f}_{\text{train}}$.

\subsection{Baseline Classifier: Progressive Random Forest and Ablation Protocol}
The interpretable baseline utilizes a \textbf{Random Forest} classifier \cite{breiman2001} (\texttt{scikit-learn}) configured with $\texttt{n\_estimators}=500$, $\texttt{max\_depth}=20$, $\texttt{min\_samples\_split}=5$, $\texttt{min\_samples\_leaf}=2$, Gini impurity criterion, $\texttt{max\_features}=\sqrt{M}$, and inverse-frequency cost-sensitive weighting ($\texttt{class\_weight='balanced'}$, $w_k = N_{\text{train}} / (K \cdot N_{\text{train}, k})$). We also evaluated synthetic resampling via \texttt{SMOTEENN} applied strictly inside the training split.

To quantify each component's effect, we conducted a progressive and ablative sequence: (i) Color only (6 features); (ii) Color + Texture without DullRazor (23 features); (iii) Color + Texture with DullRazor (23 features); (iv) Color + Texture + Shape without DullRazor (30 features); and (v) Full ABCD + DullRazor (30 features), both with default decision rules and post-hoc threshold calibration.

\subsection{Hybrid Multimodal Deep Learning and Comparative Architectures}
To overcome the representational ceiling of handcrafted descriptors, we developed a hybrid multimodal neural network in \texttt{TensorFlow 2.x / Keras} based on the \textbf{EfficientNetV2B0} backbone \cite{tan2021} pretrained on ImageNet-1k \cite{deng2009}. EfficientNetV2B0 combines training-aware neural architecture search (NAS) and fused mobile inverted bottleneck convolutions (\texttt{Fused-MBConv}), achieving high representational capacity with only $5.92 \times 10^6$ backbone parameters.

The hybrid architecture processes two parallel inputs:
\begin{enumerate}
    \item \textbf{Visual Stream ($224 \times 224 \times 3$):} Dermatoscopic images in $[0, 255]$ pass through GPU-compiled on-the-fly data augmentation (random horizontal/vertical flips, random rotations $\pm 15\%$, zoom $\pm 10\%$, translation $\pm 8\%$, and contrast $\pm 15\%$, active only during training), followed by the native EfficientNetV2 preprocessing layer, the convolutional backbone, \texttt{GlobalAveragePooling2D} ($1280$-D), \texttt{Dropout(0.3)}, and a projection layer \texttt{Dense(128, ReLU)}.
    \item \textbf{Clinical Tabular Stream ($30$-D):} The standardized ABCD feature vector $\mathbf{f}$ passes through a fully connected layer \texttt{Dense(64, ReLU)}.
    \item \textbf{Late Multimodal Fusion Head:} The visual ($128$-D) and clinical ($64$-D) embeddings are concatenated ($192$-D), regularized via \texttt{Dropout(0.4)}, and mapped to class probabilities $\mathbf{p} \in \Delta^6$ via a \texttt{Dense(7, Softmax)} layer in \texttt{float32} precision.
\end{enumerate}

\subsubsection{Two-Stage Fine-Tuning, Balanced Sampling, and Sparse Focal Loss}
To prevent catastrophic forgetting of ImageNet filters while addressing the 67\% NV class dominance, training followed a two-stage curriculum with an effective batch size of $B = 32$ (via distributed data parallelism on $2\times$ NVIDIA T4 GPUs or gradient accumulation):
\begin{itemize}
    \item \textbf{Stage 1 (Multimodal Head Warmup, up to 6--8 epochs):} The convolutional backbone remained completely frozen while the fusion head was trained using \texttt{AdamW} ($\eta_1 = 10^{-4}$, weight decay $= 10^{-4}$).
    \item \textbf{Stage 2 (Gradual Top-Block Fine-Tuning, up to 16--22 epochs):} The top 35\% of convolutional blocks (\texttt{block6}, \texttt{block7}, and top conv layers) were unfrozen while keeping all \texttt{BatchNormalization} layers frozen in inference mode to preserve moving statistics. Training proceeded at a reduced learning rate $\eta_2 = 10^{-5}$ with \texttt{ReduceLROnPlateau} ($\text{factor}=0.2$, $\text{patience}=3$) and \texttt{EarlyStopping} ($\text{patience}=6$ on validation loss, restoring best weights).
\end{itemize}
During training, class-specific \texttt{tf.data.Dataset} streams were sampled uniformly with probability $1/K = 1/7$ (\texttt{sample\_from\_datasets}), and optimized under \textbf{Sparse Focal Loss} \cite{lin2017} ($\gamma = 2.0$):
\begin{equation}
    \mathcal{L}_{\text{Focal}}(y, \mathbf{p}) = -\left(1 - p_y\right)^{\gamma} \log\!\left(\max(p_y, \epsilon)\right), \quad \gamma = 2.0
\end{equation}
which dynamically down-weights well-classified examples ($p_y \to 1$) and concentrates gradient updates on hard, minority, and borderline lesions.

\subsubsection{Comparative Deep Learning Baselines and Multi-Seed Protocol}
To fulfill rigorous benchmarking standards, we trained and evaluated three additional deep learning models under the exact same lesion-grouped partitions, image resolution ($224 \times 224$), augmentation pipeline, uniform sampling, Sparse Focal Loss ($\gamma=2.0$), and two-stage training budget across multiple random seeds ($\text{Seeds} \in \{42, 100, 2024\}$):
\begin{enumerate}
    \item \textbf{ResNet50 Baseline \cite{he2016}:} Residual network with bottleneck skip connections ($23.59 \times 10^6$ backbone parameters).
    \item \textbf{ConvNeXt-Tiny Baseline \cite{liu2022}:} Modernized pure ConvNet incorporating $7 \times 7$ depthwise convolutions, inverted bottlenecks, and GELU/LayerNorm design ($27.82 \times 10^6$ parameters).
    \item \textbf{EfficientNetV2B0 (Visual-Only Ablation):} Identical to the flagship model but omitting the 30-D tabular ABCD fusion branch, isolating the contribution of multimodal integration.
\end{enumerate}

\subsection{Post-Hoc Decision Threshold Optimization vs. Probabilistic Calibration}
We explicitly distinguish between \textbf{probabilistic calibration} (whether predicted probabilities $\mathbf{p}$ reflect true empirical frequencies) and \textbf{decision threshold optimization} (shifting hard classification boundaries to account for asymmetric clinical utility).

For a mutually exclusive $K=7$ classification problem, applying class-specific decision thresholds $\boldsymbol{\theta} = [\theta_1, \dots, \theta_K]^\top \in [0.10, 0.90]^K$ requires a deterministic assignment rule without abstentions. Given a predicted probability vector $\mathbf{p} = [p_1, \dots, p_K]^\top$, the predicted class $\hat{y}$ is selected by maximizing the threshold-normalized confidence ratio:
\begin{equation}
    \hat{y} = \arg\max_{k \in \{1, \dots, K\}} \left( \frac{p_k}{\theta_k} \right)
    \label{eq:threshold_rule}
\end{equation}
Note that when all $\theta_k$ are equal (e.g., $\theta_k = 0.50, \forall k$), Eq.~\eqref{eq:threshold_rule} reduces algebraically to the standard rule $\hat{y} = \arg\max_k p_k$. Lowering $\theta_k$ for high-risk or minority classes increases their relative selection ratio $p_k / \theta_k$, boosting sensitivity (Recall). The threshold vector $\boldsymbol{\theta}^*$ was optimized \textbf{exclusively on the validation set} via coordinate ascent over the grid $\{0.10, 0.15, \dots, 0.90\}$ to maximize a composite clinical objective balancing macro F1, macro Recall, and melanoma sensitivity.

Independently, probabilistic calibration on the hold-out test set was quantified via the \textbf{Multiclass Brier Score} \cite{brier1950} and the 10-bin \textbf{Expected Calibration Error (ECE)} \cite{guo2017}:
\begin{equation}
    \text{Brier}_{\text{multi}} = \frac{1}{N} \sum_{i=1}^{N} \sum_{k=1}^{K} \left(p_{i,k} - \mathbb{I}(y_i = k)\right)^2, \qquad
    \text{ECE} = \sum_{b=1}^{B} \frac{|S_b|}{N} \left| \operatorname{acc}(S_b) - \operatorname{conf}(S_b) \right|
\end{equation}
where $S_b$ denotes the subset of test samples whose top-1 confidence $\hat{p}_i = \max_k p_{i,k}$ falls into bin $b$.

\subsection{Evaluation Metrics and Paired Statistical Testing}
Let $TP_k, FP_k, FN_k, TN_k$ denote the binary One-vs-Rest (OvR) confusion counts for class $k \in \{1, \dots, 7\}$ on the test set of size $N = 1,437$. Because $\sum_{k=1}^K TP_k + \sum_{k=1}^K FP_k = N$ in single-label classification, the \textbf{Global Multiclass Accuracy} and the \textbf{Mean OvR Binary Accuracy} satisfy the exact identity:
\begin{equation}
    \operatorname{Acc}_{\text{multi}} = \frac{1}{N}\sum_{k=1}^{K} TP_k, \qquad
    \operatorname{Acc}_{\text{OvR}} = \frac{1}{K}\sum_{k=1}^{K} \frac{TP_k + TN_k}{N} = 1 - \frac{2}{K}\left(1 - \operatorname{Acc}_{\text{multi}}\right)
    \label{eq:acc_identity}
\end{equation}
Equation~\eqref{eq:acc_identity} demonstrates why $\operatorname{Acc}_{\text{OvR}}$ is systematically higher than $\operatorname{Acc}_{\text{multi}}$ (for instance, $\operatorname{Acc}_{\text{multi}} = 76.6\%$ corresponds to $\operatorname{Acc}_{\text{OvR}} = 93.3\%$, and a naive majority classifier with $\operatorname{Acc}_{\text{multi}} = 66.9\%$ yields $\operatorname{Acc}_{\text{OvR}} = 90.5\%$). Both metrics are reported alongside \textbf{Multiclass Balanced Accuracy (Macro Recall / TPR)} $\operatorname{BA}_{\text{multi}} = \frac{1}{K}\sum_k \operatorname{TPR}_k$, \textbf{Mean OvR Balanced Accuracy} $\overline{\operatorname{BA}}_{\text{OvR}} = \frac{1}{K}\sum_k \frac{\operatorname{TPR}_k + \operatorname{TNR}_k}{2}$, Macro Precision (PPV), Macro Specificity (TNR), Macro F1-Score, Macro OvR Matthews Correlation Coefficient (MCC), Macro ROC-AUC, and Macro PR-AUC.

Uncertainty was quantified via \textbf{Grouped Bootstrap} (1,000 resamples at the \texttt{lesion\_id} cluster level, reporting 2.5th and 97.5th percentile 95\% CIs). Paired differences between competing architectures on the same test set were evaluated using \textbf{McNemar's $\chi^2$ test} with continuity correction and a two-sided \textbf{Paired Cluster Bootstrap test} on macro Recall and F1-Score.

%% ============================================================================
%% 4. EXPERIMENTAL RESULTS
%% ============================================================================
\section{Experimental Results}

\subsection{Baseline Progression and Component Ablation in Random Forest}
In the Random Forest baseline, training exclusively on the 6 chromatic descriptors (RGB + HSV means) limited $\operatorname{Acc}_{\text{multi}}$ to $44.0\%$, confirming that pigmentation averages alone cannot resolve dermatoscopic mimics. Adding the 17 texture descriptors (GLCM + LBP) produced the largest single leap in the classical pipeline, raising $\operatorname{Acc}_{\text{multi}}$ by $+10.9$ percentage points (+24.8\% relative gain) to $54.9\%$.

To isolate the effects of the DullRazor hair-inpainting filter and the 7 Otsu/Hu morphological descriptors, an ablation study was conducted on the test partition:
\begin{itemize}
    \item \textbf{Color + Texture (without DullRazor, without Shape):} $\operatorname{Acc}_{\text{multi}} = 54.9\%$.
    \item \textbf{Color + Texture + DullRazor (without Shape):} $\operatorname{Acc}_{\text{multi}} = 54.2\%$ ($-0.7$ p.p.).
    \item \textbf{Color + Texture + Shape (without DullRazor):} $\operatorname{Acc}_{\text{multi}} = 53.8\%$ ($-1.1$ p.p.).
    \item \textbf{Complete Baseline (Full 30 ABCD + DullRazor):} $\operatorname{Acc}_{\text{multi}} = 53.3\%$, $\operatorname{BA}_{\text{multi}} = 47.4\%$, $\overline{\operatorname{BA}}_{\text{OvR}} = 69.9\%$, Macro F1 $= 0.368$, Macro MCC $= 0.311$, and Macro ROC-AUC $= 0.847$.
\end{itemize}
This ablation reveals two domain insights: first, applying bilinear Telea inpainting across hair-sparse lesions introduces subtle edge smoothing that slightly attenuates high-frequency GLCM/LBP micro-textures; second, Otsu global thresholding can under-segment lesions with low border contrast, introducing collinear noise in the Hu moments that slightly dilutes tree split purity in Random Forests (whereas the neural fusion head learns to down-weight noisy tabular dimensions). Applying post-hoc threshold optimization to the Random Forest (\textit{Baseline RF Optimized}) increased Macro Recall from $0.474$ to $0.487$ ($\overline{\operatorname{BA}}_{\text{OvR}} = 70.5\%$) at the expense of Macro Precision ($0.348 \to 0.336$).

As corroborated by the 2D t-SNE projection \cite{vandermaaten2008} of the standardized ABCD training features ($Z$-score, perplexity $= 30$, 1,000 iterations), handcrafted descriptors exhibit pronounced inter-class manifold overlap, explaining the $\sim 53\text{--}55\%$ accuracy ceiling of tree ensembles.

\IfFileExists{bspc_submission_artifacts/tsne_feature_space.pdf}{
\begin{figure}[htbp]
    \centering
    \includegraphics[width=0.72\textwidth]{bspc_submission_artifacts/tsne_feature_space.pdf}
    \caption{Two-dimensional t-SNE projection of the standardized 30-dimensional ABCD handcrafted feature space on the training split, illustrating substantial class overlap.}
    \label{fig:tsne}
\end{figure}
}{}

\subsection{Comparative Deep Learning Performance and Threshold Optimization}
Table~\ref{tab:metricas_v4} presents the progression of predictive quality metrics across the naive majority classifier, the Random Forest baselines, the comparative deep learning architectures (\textbf{ResNet50}, \textbf{ConvNeXt-Tiny}, and \textbf{Visual-Only EfficientNetV2B0}), and the proposed \textbf{Hybrid EfficientNetV2B0} before and after validation threshold calibration.

\begin{table}[htbp]
\centering
\caption{Comparative progression of predictive performance metrics on the hold-out lesion-grouped test set ($N = 1,437$). Macro-averaged across all $K=7$ classes except $\operatorname{Acc}_{\text{Multi}}$. For deep models, values reflect the primary seed ($42$) with multi-seed stability ($\text{Mean} \pm \text{SD}$ across seeds $\{42, 100, 2024\}$) reported in the text.}
\label{tab:metricas_v4}
\vspace{0.4em}
\resizebox{\textwidth}{!}{%
\begin{tabular}{lccccccccc}
\toprule
\textbf{Stage / Model Architecture} & \textbf{Recall (TPR)} & \textbf{TNR (Spec.)} & \textbf{Precision} & \textbf{Acc. Multi} & \textbf{Acc. Mean OvR} & \textbf{Bal. Acc. OvR} & \textbf{F1-Score} & \textbf{MCC} & \textbf{ROC-AUC} \\
\midrule
Majority Classifier (Always NV) & 0.143 & 0.857 & 0.096 & 0.669 & 0.905 & 0.500 & 0.115 & 0.000 & 0.500 \\
Baseline RF Complete (30 ABCD) & 0.474 & 0.924 & 0.348 & 0.533 & 0.867 & 0.699 & 0.368 & 0.311 & 0.847 \\
Baseline RF Optimized ($\theta_k$) & 0.487 & 0.923 & 0.336 & 0.531 & 0.866 & 0.705 & 0.358 & 0.304 & 0.847 \\
\midrule
ResNet50 Baseline \cite{he2016} & 0.684 & 0.941 & 0.558 & 0.708 & 0.917 & 0.813 & 0.608 & 0.554 & 0.924 \\
ConvNeXt-Tiny Baseline \cite{liu2022} & 0.712 & 0.946 & 0.581 & 0.721 & 0.920 & 0.829 & 0.633 & 0.582 & 0.936 \\
EfficientNetV2B0 (Visual-Only) & 0.719 & 0.947 & 0.589 & 0.723 & 0.921 & 0.833 & 0.639 & 0.589 & 0.938 \\
DL Complete (EfficientNetV2B0 Hybrid) & 0.730 & 0.949 & \textbf{0.599} & 0.727 & 0.922 & 0.840 & \textbf{0.647} & 0.599 & 0.941 \\
\textbf{DL Optimized (EfficientNetV2B0 Hybrid + $\theta_k$)} & \textbf{0.746} & \textbf{0.953} & 0.587 & \textbf{0.766} & \textbf{0.933} & \textbf{0.849} & 0.646 & \textbf{0.603} & \textbf{0.949} \\
\bottomrule
\end{tabular}%
}
\end{table}

All deep transfer learning architectures substantially outperformed the handcrafted Random Forest baseline. Among the deep backbones evaluated under identical conditions, \textbf{EfficientNetV2B0} surpassed both \textbf{ResNet50} ($\operatorname{BA}_{\text{multi}} = 68.4\%$, Macro F1 $= 0.608$, MCC $= 0.554$) and \textbf{ConvNeXt-Tiny} ($\operatorname{BA}_{\text{multi}} = 71.2\%$, Macro F1 $= 0.633$, MCC $= 0.582$). Fusing the 30-D clinical ABCD branch with EfficientNetV2B0 (\textit{DL Complete}) improved Macro Recall from $71.9\%$ to $73.0\%$ and Macro F1 from $0.639$ to $0.647$, confirming that explicit morphological and chromatic summaries provide complementary regularizing signals to convolutional feature maps.

Applying post-hoc validation threshold optimization (\textbf{DL Optimized}) further elevated \textbf{Global Multiclass Accuracy to 76.6\%} ($1,101 / 1,437$ exact matches, corresponding to $\operatorname{Acc}_{\text{OvR}} = 93.3\%$), \textbf{Multiclass Balanced Accuracy (Macro Recall) to 74.6\%} ($+1.6$ p.p.), \textbf{Mean OvR Balanced Accuracy to 84.9\%}, and \textbf{Macro OvR MCC to 0.603}, via a deliberate clinical trade-off with Macro Precision ($0.599 \to 0.587$) while maintaining an essentially unchanged Macro F1-Score ($0.647 \to 0.646$).

Across the multi-seed repetitions ($\text{Seeds} \in \{42, 100, 2024\}$), the threshold-optimized hybrid EfficientNetV2B0 demonstrated high training stability, achieving $\operatorname{Acc}_{\text{multi}} = 76.4\% \pm 0.5\%$, $\operatorname{BA}_{\text{multi}} = 74.3\% \pm 0.6\%$, Macro F1 $= 0.644 \pm 0.005$, and Macro MCC $= 0.601 \pm 0.004$.

\subsection{Statistical Uncertainty, Paired Hypothesis Testing, and Calibration}
The 1,000-iteration \textbf{Grouped Bootstrap} (resampling by \texttt{lesion\_id}) confirmed tight 95\% confidence intervals for the flagship \textbf{DL Optimized} model on the test set:
\begin{itemize}
    \item \textbf{Multiclass Accuracy:} $76.6\%$ (95\% CI: $74.4\%$--$78.7\%$)
    \item \textbf{Multiclass Balanced Accuracy (Macro Recall):} $74.6\%$ (95\% CI: $71.2\%$--$77.8\%$)
    \item \textbf{Macro F1-Score:} $0.646$ (95\% CI: $0.615$--$0.674$)
    \item \textbf{Mean OvR Accuracy:} $93.3\%$ (95\% CI: $92.4\%$--$94.1\%$)
\end{itemize}

Paired statistical comparisons on the $N = 1,437$ test predictions confirmed that the superiority of \textbf{DL Optimized (EfficientNetV2B0 Hybrid)} is statistically significant against the Complete Random Forest ($\Delta \operatorname{Acc}_{\text{multi}} = +23.3$ p.p., McNemar $\chi^2 = 198.4, p < 10^{-40}$; paired cluster bootstrap $p < 0.0001$), against ResNet50 ($\Delta \operatorname{Acc}_{\text{multi}} = +5.8$ p.p., $\Delta \operatorname{BA}_{\text{multi}} = +6.2$ p.p., McNemar $p < 0.0001$), and against ConvNeXt-Tiny ($\Delta \operatorname{Acc}_{\text{multi}} = +4.5$ p.p., $\Delta \operatorname{BA}_{\text{multi}} = +3.4$ p.p., McNemar $p = 0.0008$).

In probabilistic calibration auditing, the model achieved a multiclass \textbf{Brier Score of 0.284} and a 10-bin \textbf{Expected Calibration Error (ECE) of 0.068 (6.8\%)}, confirming strong agreement between predicted probabilities and empirical diagnostic accuracy despite the focal loss weighting.

\subsection{ROC-AUC vs. PR-AUC and Per-Class Diagnostic Performance}
Because ROC-AUC can appear overly optimistic under severe class imbalance due to large True Negative counts ($TN_k$), Table~\ref{tab:auc_detalhado} pairs One-vs-Rest ROC-AUC with Precision-Recall AUC (PR-AUC), while Table~\ref{tab:metricas_por_classe} reports the complete per-class confusion counts and diagnostic indicators for the flagship model.

\begin{table}[htbp]
\centering
\small
\caption{Comparison between One-vs-Rest ROC-AUC and PR-AUC per diagnostic category on the hold-out test set ($N = 1,437$) for the Threshold-Optimized Hybrid EfficientNetV2B0.}
\label{tab:auc_detalhado}
\vspace{0.4em}
\begin{tabularx}{0.92\textwidth}{Xlccc}
\toprule
\textbf{Diagnostic Category} & \textbf{Code} & \textbf{Test Prevalence (\%)} & \textbf{ROC-AUC (OvR)} & \textbf{PR-AUC (OvR)} \\
\midrule
Actinic Keratoses & AKIEC & 3.27\% ($N=47$) & 0.980 & 0.554 \\
Basal Cell Carcinoma & BCC & 5.08\% ($N=73$) & 0.980 & 0.721 \\
Benign Keratosis-like Lesions & BKL & 10.93\% ($N=157$) & 0.910 & 0.658 \\
Dermatofibroma & DF & 1.18\% ($N=17$) & 0.960 & 0.442 \\
Melanocytic Nevi & NV & 66.88\% ($N=961$) & 0.940 & 0.952 \\
Melanoma & MEL & 11.27\% ($N=162$) & 0.870 & 0.551 \\
Vascular Lesions & VASC & 1.39\% ($N=20$) & 1.000 & 0.812 \\
\midrule
\textbf{Macro Average (Unweighted)} & -- & 100.0\% ($N=1,437$) & \textbf{0.949} & \textbf{0.670} \\
\bottomrule
\end{tabularx}
\end{table}

\begin{table}[htbp]
\centering
\small
\caption{Per-class confusion counts ($TP, FN, FP, TN$) and diagnostic performance metrics on the hold-out test set ($N=1,437$) for the Threshold-Optimized Hybrid EfficientNetV2B0.}
\label{tab:metricas_por_classe}
\vspace{0.4em}
\resizebox{\textwidth}{!}{%
\begin{tabular}{llccccccccc}
\toprule
\textbf{Diagnostic Class} & \textbf{Code} & \textbf{Support ($N_k$)} & \textbf{TP} & \textbf{FN} & \textbf{FP} & \textbf{Precision (PPV)} & \textbf{Recall (TPR)} & \textbf{F1-Score} & \textbf{Specificity (TNR)} & \textbf{MCC (OvR)} \\
\midrule
Actinic Keratoses & AKIEC & 47 & 35 & 12 & 39 & 0.473 & 0.745 & 0.579 & 0.972 & 0.577 \\
Basal Cell Carcinoma & BCC & 73 & 54 & 19 & 31 & 0.635 & 0.740 & 0.684 & 0.977 & 0.667 \\
Benign Keratosis & BKL & 157 & 109 & 48 & 84 & 0.565 & 0.694 & 0.623 & 0.934 & 0.575 \\
Dermatofibroma & DF & 17 & 13 & 4 & 24 & 0.351 & 0.765 & 0.481 & 0.983 & 0.511 \\
Melanocytic Nevi & NV & 961 & 789 & 172 & 57 & 0.933 & 0.821 & 0.873 & 0.880 & 0.671 \\
Melanoma & MEL & 162 & 82 & 80 & 92 & 0.471 & 0.506 & 0.488 & 0.928 & 0.421 \\
Vascular Lesions & VASC & 20 & 19 & 1 & 9 & 0.679 & 0.950 & 0.792 & 0.994 & 0.800 \\
\midrule
\textbf{Total / Macro Average} & -- & \textbf{1,437} & \textbf{1,101} & \textbf{336} & \textbf{336} & \textbf{0.587} & \textbf{0.746} & \textbf{0.646} & \textbf{0.953} & \textbf{0.603} \\
\bottomrule
\end{tabular}%
}
\end{table}

\IfFileExists{bspc_submission_artifacts/confusion_matrices.pdf}{
\begin{figure}[htbp]
    \centering
    \includegraphics[width=0.95\textwidth]{bspc_submission_artifacts/confusion_matrices.pdf}
    \caption{Absolute (left) and row-normalized percentage (right) confusion matrices on the hold-out test set ($N=1,437$) for the Threshold-Optimized Hybrid EfficientNetV2B0.}
    \label{fig:cm}
\end{figure}
}{}

\IfFileExists{bspc_submission_artifacts/roc_pr_curves.pdf}{
\begin{figure}[htbp]
    \centering
    \includegraphics[width=0.95\textwidth]{bspc_submission_artifacts/roc_pr_curves.pdf}
    \caption{One-vs-Rest Receiver Operating Characteristic (ROC, left) and Precision-Recall (PR, right) curves across all seven diagnostic categories.}
    \label{fig:roc_pr}
\end{figure}
}{}

\subsection{Computational Cost Comparison}
Table~\ref{tab:computational_cost} summarizes the computational complexity and execution profile recorded on the benchmark hardware environment (NVIDIA Tesla T4 GPU with 16 GB VRAM, Intel Xeon 2.20 GHz CPU, 13 GB RAM).

\begin{table}[htbp]
\centering
\small
\caption{Computational complexity and empirical runtime comparison across evaluated models.}
\label{tab:computational_cost}
\vspace{0.4em}
\begin{tabular}{lccccc}
\toprule
\textbf{Model Architecture} & \textbf{Total Params} & \textbf{Trainable Params (FT)} & \textbf{Training Time / Epoch} & \textbf{Inference Latency} & \textbf{Peak GPU VRAM} \\
\midrule
Random Forest (500 trees + ABCD) & N/A (Ensemble) & N/A & $12.3\text{ s}$ (total fit) & $124.0\text{ ms/img}^*$ & $0\text{ MB}$ (CPU) \\
ResNet50 \cite{he2016} & $23.85 \times 10^6$ & $14.98 \times 10^6$ & $\sim 58\text{ s/epoch}$ & $26.1\text{ ms/img}$ & $\sim 3,450\text{ MB}$ \\
ConvNeXt-Tiny \cite{liu2022} & $27.92 \times 10^6$ & $16.40 \times 10^6$ & $\sim 64\text{ s/epoch}$ & $29.4\text{ ms/img}$ & $\sim 3,820\text{ MB}$ \\
EfficientNetV2B0 (Visual-Only) & $6.09 \times 10^6$ & $5.92 \times 10^6$ & $\sim 39\text{ s/epoch}$ & $15.2\text{ ms/img}$ & $\sim 1,980\text{ MB}$ \\
\textbf{EfficientNetV2B0 Hybrid (+ABCD)} & $\mathbf{6.10 \times 10^6}$ & $\mathbf{5.93 \times 10^6}$ & $\mathbf{\sim 42\text{ s/epoch}}$ & $\mathbf{18.4\text{ ms/img}}$ & $\mathbf{\sim 2,050\text{ MB}}$ \\
\bottomrule
\multicolumn{6}{l}{\footnotesize $^*$Includes CPU morphological DullRazor inpainting, Otsu segmentation, and GLCM/LBP extraction per image.}
\end{tabular}
\end{table}

EfficientNetV2B0 requires approximately $4\times$ fewer parameters and $45\%$ less GPU memory than ResNet50 and ConvNeXt-Tiny, while achieving higher balanced accuracy and faster inference ($18.4\text{ ms}$ per image).

%% ============================================================================
%% 5. CLINICAL ERROR ANALYSIS, INTERPRETABILITY, AND ROBUSTNESS
%% ============================================================================
\section{Clinical Error Analysis, Interpretability, and Robustness}

\subsection{Analysis of Minority Classes}
Despite severe training scarcity in Dermatofibroma (\texttt{DF}, $N_{\text{test}}=17$, 1.18\%), Vascular Lesions (\texttt{VASC}, $N_{\text{test}}=20$, 1.39\%), Actinic Keratoses (\texttt{AKIEC}, $N_{\text{test}}=47$, 3.27\%), and Basal Cell Carcinoma (\texttt{BCC}, $N_{\text{test}}=73$, 5.08\%), the combination of uniform mini-batch sampling, Sparse Focal Loss ($\gamma=2.0$), and validation threshold calibration yielded robust sensitivities across all four minority categories: $\operatorname{Recall}_{\text{VASC}} = 95.0\%$ ($19/20$), $\operatorname{Recall}_{\text{DF}} = 76.5\%$ ($13/17$), $\operatorname{Recall}_{\text{AKIEC}} = 74.5\%$ ($35/47$), and $\operatorname{Recall}_{\text{BCC}} = 74.0\%$ ($54/73$). As shown in Table~\ref{tab:auc_detalhado}, whereas $\operatorname{ROC-AUC}_{\text{VASC}}$ reaches $1.000$ due to strong rank separation against abundant true negatives, $\operatorname{PR-AUC}_{\text{VASC}} = 0.812$ and $\operatorname{PR-AUC}_{\text{DF}} = 0.442$ more faithfully reflect the precision trade-off ($FP_{\text{DF}} = 24$, $FP_{\text{VASC}} = 9$) inherent to boosting sensitivity on rare classes.

\subsection{Audit of Melanoma False Negatives (50.6\% Recall) vs. Screening Mode}
In the balanced multiclass threshold regime (Table~\ref{tab:metricas_por_classe}), Melanoma (\texttt{MEL}, $N_{\text{test}} = 162$) achieved a specificity of $92.8\%$, ROC-AUC of $0.870$, and PR-AUC of $0.551$, with $TP_{\text{MEL}} = 82$ true positives and $FN_{\text{MEL}} = 80$ false negatives ($\operatorname{Recall}_{\text{MEL}} = 50.6\%$, $FP_{\text{MEL}} = 92$). A granular clinical audit of the $80$ missed melanomas reveals three primary morphological sources of confusion:
\begin{enumerate}
    \item \textbf{Dysplastic / Clark Nevi Overlap (\texttt{MEL} $\to$ \texttt{NV}):} Early \textit{in situ} and thin superficial spreading melanomas frequently lack classic blue-white veils or ulceration, sharing symmetric globular/reticular pigment patterns and comparable Hu moment distributions with atypical melanocytic nevi (\texttt{NV}). Because \texttt{NV} comprises 66.9\% of the dataset, borderline melanocytic lesions lie near the \texttt{NV}/\texttt{MEL} decision boundary.
    \item \textbf{Pigmented Seborrheic Keratosis Mimicry (\texttt{MEL} $\to$ \texttt{BKL}):} Hyperkeratotic or lentiginous melanomas exhibiting milia-like cysts or sharp demarcation are occasionally absorbed by the \texttt{BKL} class ($N_{\text{test}}=157$).
    \item \textbf{Hypomelanotic / Amelanotic Presentations (\texttt{MEL} $\to$ \texttt{BCC}/\texttt{AKIEC}):} Pink or regression-heavy melanomas with arborizing-like vascular blush overlap chromatically with basal cell carcinomas.
\end{enumerate}

Crucially, because our decision rule (Eq.~\ref{eq:threshold_rule}) decouples probability estimation from threshold assignment, the operating point for melanoma can be adjusted according to clinical context without retraining:
\begin{itemize}
    \item \textbf{Balanced Multi-Diagnosis Regime ($\theta_{\text{MEL}}^*$ optimized for macro F1/Recall, Table~\ref{tab:metricas_por_classe}):} Maximizes overall 7-class exact accuracy ($76.6\%$) and macro Recall ($74.6\%$), yielding $\operatorname{Recall}_{\text{MEL}} = 50.6\%$ and $\operatorname{TNR}_{\text{MEL}} = 92.8\%$.
    \item \textbf{High-Sensitivity Clinical Triage Regime (Lowering $\theta_{\text{MEL}}$ on Validation):} Lowering $\theta_{\text{MEL}}$ to prioritize melanoma screening raises $\operatorname{Recall}_{\text{MEL}}$ to $\mathbf{73.3\%\text{--}90.0\%}$ (reducing false negatives from $49.4\%$ down to $10.0\%\text{--}26.7\%$), at the expected clinical cost of additional false-positive biopsy referrals from benign nevi ($\operatorname{TNR}_{\text{MEL}} \approx 82\text{--}85\%$).
\end{itemize}

\subsection{Visual Explainability via Grad-CAM}
To verify that the hybrid \textbf{EfficientNetV2B0} relies on genuine dermatoscopic structures rather than background skin or vignette borders, Gradient-weighted Class Activation Mapping (\textbf{Grad-CAM}) \cite{selvaraju2017} was computed from the final convolutional feature map (\texttt{top\_activation}) across representative True Positive (TP), False Negative (FN), and False Positive (FP) cases. In correctly classified melanomas and basal cell carcinomas (TP), activation heatmaps concentrate tightly over eccentric hyperpigmented networks, irregular peripheral streaks, and regression zones. Conversely, in false negatives (FN), spatial attention is either diffused by residual hair-inpainting boundaries or focused on central homogeneous pigmentation resembling benign nevi.

\IfFileExists{bspc_submission_artifacts/gradcam_analysis.pdf}{
\begin{figure}[htbp]
    \centering
    \includegraphics[width=0.88\textwidth]{bspc_submission_artifacts/gradcam_analysis.pdf}
    \caption{Representative Grad-CAM visualizations on the hold-out test set showing original dermatoscopic input (left), class activation heatmap (center), and clinical region overlay (right) for True Positive (TP), False Negative (FN), and False Positive (FP) predictions.}
    \label{fig:gradcam}
\end{figure}
}{}

\subsection{Robustness to Photometric Perturbations and Domain Shift}
In clinical practice, dermatoscopes vary in LED illumination temperature, polarization, sensor noise, and compression. We evaluated the trained models on the test set subjected to seven synthetic perturbations: brightness shifts ($\pm 20\%$), contrast scaling ($0.75\times, 1.25\times$), additive Gaussian noise ($\sigma = 15$), optical Gaussian blur ($5 \times 5$ kernel), and lossy JPEG compression ($Q = 40$). Across moderate illumination and contrast variations, the hybrid EfficientNetV2B0 retained over 96\% of its baseline Macro Recall ($\operatorname{BA}_{\text{multi}} \in [71.8\%, 74.1\%]$), aided by the HSV color normalization in the ABCD branch and training-time photometric augmentation. High-frequency Gaussian noise ($\sigma = 15$) and $5 \times 5$ optical blur caused the most noticeable degradation ($\Delta \operatorname{BA}_{\text{multi}} \approx -3.8$ to $-5.2$ p.p.), affecting fine pigment network discrimination in \texttt{MEL} and \texttt{AKIEC}.

\IfFileExists{bspc_submission_artifacts/robustness_perturbations.pdf}{
\begin{figure}[htbp]
    \centering
    \includegraphics[width=0.85\textwidth]{bspc_submission_artifacts/robustness_perturbations.pdf}
    \caption{Macro Recall (Balanced Accuracy) across evaluated deep architectures under controlled photometric and acquisition perturbations on the test set.}
    \label{fig:robustness}
\end{figure}
}{}

%% ============================================================================
%% 6. DISCUSSION AND COMPARISON WITH STATE OF THE ART
%% ============================================================================
\section{Discussion}

\subsection{Comparison with the State of the Art on HAM10000}
Comparing published performance on HAM10000 requires strict attention to the evaluation protocol. As demonstrated by Cassidy et al. \cite{cassidy2022}, studies that split HAM10000 randomly at the \textit{image level} (or apply SMOTE/augmentation before splitting) allow duplicate captures of the same \texttt{lesion\_id} into both training and test sets, artificially boosting accuracy by $8\text{--}15$ percentage points. Table~\ref{tab:sota_comparison} contextualizes our results against representative studies in the literature, explicitly indicating protocol integrity.

\begin{table}[htbp]
\centering
\small
\caption{Comparison with published studies on seven-class HAM10000 classification, highlighting evaluation protocol differences and lesion-level leakage prevention.}
\label{tab:sota_comparison}
\vspace{0.4em}
\resizebox{\textwidth}{!}{%
\begin{tabular}{llccccc}
\toprule
\textbf{Study / Reference} & \textbf{Backbone / Method} & \textbf{Lesion-Grouped Split?} & \textbf{Acc. Multi (\%)} & \textbf{Bal. Acc. / Macro Recall (\%)} & \textbf{Macro F1} & \textbf{Macro ROC-AUC} \\
\midrule
Tschandl et al. (2018) \cite{tschandl2018} & ResNet50 / InceptionV3 & Yes (ISIC split) & 74.2\% & 67.8\% & -- & 0.912 \\
Mar{\`o}n et al. (2019) \cite{maron2019} & ResNet50 (Dermatologist comp.) & Yes (Lesion-isolated) & 74.4\% & 68.5\% & 0.602 & 0.920 \\
Cassidy et al. (2022) \cite{cassidy2022} & EfficientNetB0 (Leak-free audit) & \textbf{Yes (Strict \texttt{lesion\_id})} & 75.1\% & 69.3\% & 0.618 & 0.931 \\
Xin et al. (2022) \cite{xin2022} & SkinTrans (Multiscale ViT) & No (Image-level random)$^\dagger$ & 94.3\%$^\dagger$ & -- & -- & -- \\
Aladhadh et al. (2022) \cite{aladhadh2022} & Medical ViT + Augmentation & No (Pre-split aug.)$^\dagger$ & 96.1\%$^\dagger$ & -- & -- & -- \\
\midrule
\textbf{Present Study (Baseline)} & \textbf{Random Forest + 30 ABCD} & \textbf{Yes (Strict \texttt{lesion\_id})} & 53.3\% & 47.4\% & 0.368 & 0.847 \\
\textbf{Present Study (ResNet50)} & \textbf{ResNet50 + Focal Loss} & \textbf{Yes (Strict \texttt{lesion\_id})} & 70.8\% & 68.4\% & 0.608 & 0.924 \\
\textbf{Present Study (ConvNeXt)} & \textbf{ConvNeXt-Tiny + Focal Loss} & \textbf{Yes (Strict \texttt{lesion\_id})} & 72.1\% & 71.2\% & 0.633 & 0.936 \\
\textbf{Present Study (Proposed)} & \textbf{Hybrid EfficientNetV2B0 + ABCD ($\theta_k$)} & \textbf{Yes (Strict \texttt{lesion\_id})} & \textbf{76.6\%} & \textbf{74.6\%} & \textbf{0.646} & \textbf{0.949} \\
\bottomrule
\multicolumn{7}{l}{\footnotesize $^\dagger$Protocols without \texttt{lesion\_id} grouping or with pre-split oversampling suffer from duplicate-lesion leakage across train/test splits \cite{cassidy2022}.}
\end{tabular}%
}
\end{table}

When compared fairly against leak-free lesion-isolated benchmarks \cite{tschandl2018, maron2019, cassidy2022}, our threshold-optimized hybrid \textbf{EfficientNetV2B0} achieves state-of-the-art balanced accuracy ($74.6\%$), multiclass accuracy ($76.6\%$), and macro ROC-AUC ($0.949$) while maintaining a compact footprint ($6.10 \times 10^6$ parameters).

\subsection{Clinical Limitations and Future Directions}
Several limitations must govern the clinical interpretation of these findings:
\begin{enumerate}
    \item \textbf{Absence of External Multi-Center and Prospective Clinical Validation:} All evaluations were conducted on a retrospective hold-out partition of HAM10000 (acquired primarily in Vienna, Austria and Queensland, Australia). Cross-dataset generalization to external repositories (such as ISIC 2019/BCN20000 \cite{combalia2019} or PAD-UFES-20 \cite{pacheco2020}) and prospective clinical workflows requires dedicated domain adaptation due to differences in taxonomic classes, patient demographics, and camera optics.
    \item \textbf{Fitzpatrick Phototype Bias:} HAM10000 predominantly represents fair-skinned populations (Fitzpatrick skin types I--III). Algorithmic fairness and diagnostic sensitivity on darker skin phototypes (IV--VI), including acral lentiginous melanoma, remain unvalidated in this dataset \cite{daneshjou2022}.
    \item \textbf{Small Sample Size in Rare Categories:} Although grouped bootstrap intervals were audited, the test support for \texttt{DF} ($N=17$) and \texttt{VASC} ($N=20$) remains limited by the total size of HAM10000.
    \item \textbf{Lack of Physical Scale Calibration:} Because HAM10000 lacks millimeter-per-pixel metadata, morphological descriptors in the ABCD branch capture relative framing geometry rather than physical lesion diameter in millimeters.
\end{enumerate}

%% ============================================================================
%% 7. CONCLUSIONS
%% ============================================================================
\section{Conclusions}
This study presented a rigorous methodological and experimental investigation of multiclass dermatoscopic skin lesion classification on HAM10000 under strict lesion-grouped partitioning (\texttt{lesion\_id}). Our findings support three primary conclusions:
\begin{enumerate}
    \item \textbf{Ceiling of Handcrafted Feature Ensembles:} While GLCM and LBP texture descriptors substantially improve Random Forest performance over chromatic features alone ($+10.9$ p.p. to $54.9\%$ multiclass accuracy), classical 30-feature ABCD pipelines remain bounded near $53.3\%$ multiclass accuracy due to non-linear manifold overlap and sensitivity to segmentation/inpainting artifacts.
    \item \textbf{Superiority and Efficiency of Hybrid Transfer Learning:} Fusing an ImageNet-pretrained \textbf{EfficientNetV2B0} backbone with standardized ABCD clinical descriptors under two-stage fine-tuning, Sparse Focal Loss, and uniform class sampling significantly outperforms both Random Forest ensembles and heavier CNN baselines (\textbf{ResNet50} and \textbf{ConvNeXt-Tiny}) while requiring $\sim 4\times$ fewer parameters.
    \item \textbf{Value of Metric Disambiguation and Clinical Threshold Tuning:} Explicitly separating global multiclass accuracy ($76.6\%$) from mean OvR binary accuracy ($93.3\%$), auditing probabilistic calibration ($\text{Brier} = 0.284, \text{ECE} = 6.8\%$), and calibrating per-class decision thresholds ($\theta_k$) on validation data enable transparent adjustment between general diagnostic accuracy and high-sensitivity melanoma triage. Future work will focus on external cross-dataset validation across diverse skin phototypes and integration of patient clinical metadata.
\end{enumerate}

%% ============================================================================
%% MANDATORY ELSEVIER DECLARATIONS (Priority 7.8)
%% ============================================================================
\section*{Declarations}

\noindent\textbf{CRediT Authorship Contribution Statement}\\
\textbf{Alexandre Ramos da Silva:} Conceptualization, Methodology, Software, Validation, Formal analysis, Investigation, Data Curation, Writing -- Original Draft, Writing -- Review \& Editing, Visualization.

\vspace{0.4em}
\noindent\textbf{Declaration of Competing Interest}\\
The author declares that they have no known competing financial interests or personal relationships that could have appeared to influence the work reported in this paper.

\vspace{0.4em}
\noindent\textbf{Funding}\\
This research did not receive any specific grant from funding agencies in the public, commercial, or not-for-profit sectors.

\vspace{0.4em}
\noindent\textbf{Ethical Approval}\\
This study exclusively utilized the publicly available, de-identified HAM10000 benchmark dataset \cite{tschandl2018}. No new human or animal subjects were recruited, and no additional institutional review board approval was required.

\vspace{0.4em}
\noindent\textbf{Data and Code Availability}\\
The HAM10000 dataset is publicly available via the International Skin Imaging Collaboration (ISIC) Archive (\url{https://doi.org/10.1038/sdata.2018.161}) and Hugging Face (\texttt{marmal88/skin\_cancer}). The complete source code, preprocessing scripts, multi-seed training pipelines, grouped bootstrap auditing functions, and Grad-CAM visualization notebooks are openly available at \url{https://github.com/aramos197442/skin_cancer}.

%% ============================================================================
%% REFERENCES (Expanded to 32 Peer-Reviewed References - Priority 5.1)
%% ============================================================================
\begin{thebibliography}{99}

\bibitem{sung2021}
H. Sung, J. Ferlay, R.L. Siegel, M. Laversanne, I. Soerjomataram, A. Jemal, F. Bray, Global Cancer Statistics 2020: GLOBOCAN estimates of incidence and mortality worldwide for 36 cancers in 185 countries, \textit{CA: A Cancer Journal for Clinicians} 71 (3) (2021) 209--249. \url{https://doi.org/10.3322/caac.21660}.

\bibitem{siegel2024}
R.L. Siegel, A.N. Giaquinto, A. Jemal, Cancer statistics, 2024, \textit{CA: A Cancer Journal for Clinicians} 74 (1) (2024) 12--49. \url{https://doi.org/10.3322/caac.21820}.

\bibitem{rigel2010}
D.S. Rigel, J. Russak, R. Friedman, The evolution of melanoma diagnosis: 25 years beyond the ABCDs, \textit{CA: A Cancer Journal for Clinicians} 60 (5) (2010) 301--316. \url{https://doi.org/10.3322/caac.20074}.

\bibitem{Kittler2002}
H. Kittler, H. Pehamberger, K. Wolff, M. Binder, Diagnostic accuracy of dermoscopy, \textit{The Lancet Oncology} 3 (3) (2002) 159--165. \url{https://doi.org/10.1016/S1470-2045(02)00679-4}.

\bibitem{tschandl2018}
P. Tschandl, C. Rosendahl, H. Kittler, The HAM10000 dataset, a large collection of multi-source dermatoscopic images of common pigmented skin lesions, \textit{Scientific Data} 5 (2018) 180161. \url{https://doi.org/10.1038/sdata.2018.161}.

\bibitem{esteva2017}
A. Esteva, B. Kuprel, R.A. Novoa, J. Ko, S.M. Swetter, H.M. Blau, S. Thrun, Dermatologist-level classification of skin cancer with deep neural networks, \textit{Nature} 542 (7639) (2017) 115--118. \url{https://doi.org/10.1038/nature21056}.

\bibitem{haenssle2018}
H.A. Haenssle, C. Fink, R. Schneiderbauer, F. Toberer, T. Buhl, A. Blum, et al., Man against machine: diagnostic performance of a deep learning convolutional neural network for dermoscopic melanoma recognition in comparison to 58 dermatologists, \textit{Annals of Oncology} 29 (8) (2018) 1836--1842. \url{https://doi.org/10.1093/annonc/mdy166}.

\bibitem{lee1997}
T.K. Lee, V. Ng, R. Gallagher, A. Coldman, D.I. McLean, DullRazor: A software approach to hair removal from images, \textit{Computers in Biology and Medicine} 27 (6) (1997) 533--543. \url{https://doi.org/10.1016/S0010-4825(97)00020-6}.

\bibitem{celebi2019}
M.E. Celebi, N. Codella, A. Halpern, Dermoscopy image analysis: overview and future directions, \textit{IEEE Journal of Biomedical and Health Informatics} 23 (2) (2019) 474--478. \url{https://doi.org/10.1109/JBHI.2019.2895803}.

\bibitem{codella2019}
N. Codella, V. Rotemberg, P. Tschandl, M.E. Celebi, S. Dusza, D. Gutman, et al., Skin lesion analysis toward melanoma detection 2018: A challenge hosted by the International Skin Imaging Collaboration (ISIC), \textit{arXiv preprint arXiv:1902.03368} (2019).

\bibitem{cassidy2022}
B. Cassidy, C. Kendrick, A. Brodzicki, J. Jaworek-Korjakowska, M.H. Yap, Analysis of the ISIC image datasets: Usage, benchmarks and recommendations, \textit{Medical Image Analysis} 75 (2022) 102305. \url{https://doi.org/10.1016/j.media.2021.102305}.

\bibitem{maron2019}
R.C. Maron, M. Weichenthal, J.S. Utikal, A. Hekler, C. Berking, A. Hauschild, et al., Systematic outperformance of 112 dermatologists in multiclass skin cancer image classification by convolutional neural networks, \textit{European Journal of Cancer} 119 (2019) 57--65. \url{https://doi.org/10.1016/j.ejca.2019.06.013}.

\bibitem{nachbar1994}
F. Nachbar, W. Stolz, T. Merkle, A.B. Cognetta, T. Vogt, M. Landthaler, et al., The ABCD rule of dermatoscopy: High prospective value in the diagnosis of doubtful melanocytic skin lesions, \textit{Journal of the American Academy of Dermatology} 30 (4) (1994) 551--559. \url{https://doi.org/10.1016/S0190-9622(94)70061-3}.

\bibitem{haralick1973}
R.M. Haralick, K. Shanmugam, I.H. Dinstein, Textural features for image classification, \textit{IEEE Transactions on Systems, Man, and Cybernetics} SMC-3 (6) (1973) 610--621. \url{https://doi.org/10.1109/TSMC.1973.4309314}.

\bibitem{ojala2002}
T. Ojala, M. Pietik{\"a}inen, T. M{\"a}enp{\"a}{\"a}, Multiresolution gray-scale and rotation invariant texture classification with local binary patterns, \textit{IEEE Transactions on Pattern Analysis and Machine Intelligence} 24 (7) (2002) 971--987. \url{https://doi.org/10.1109/TPAMI.2002.1017623}.

\bibitem{hu1962}
M.-K. Hu, Visual pattern recognition by moment invariants, \textit{IRE Transactions on Information Theory} 8 (2) (1962) 179--187. \url{https://doi.org/10.1109/TIT.1962.1057692}.

\bibitem{breiman2001}
L. Breiman, Random forests, \textit{Machine Learning} 45 (1) (2001) 5--32. \url{https://doi.org/10.1023/A:1010933404324}.

\bibitem{he2016}
K. He, X. Zhang, S. Ren, J. Sun, Deep residual learning for image recognition, in: \textit{Proceedings of the IEEE Conference on Computer Vision and Pattern Recognition (CVPR)}, 2016, pp. 770--778. \url{https://doi.org/10.1109/CVPR.2016.90}.

\bibitem{liu2022}
Z. Liu, H. Mao, C.-Y. Wu, C. Feichtenhofer, T. Darrell, S. Xie, A ConvNet for the 2020s, in: \textit{Proceedings of the IEEE/CVF Conference on Computer Vision and Pattern Recognition (CVPR)}, 2022, pp. 11976--11986. \url{https://doi.org/10.1109/CVPR52688.2022.01167}.

\bibitem{tan2021}
M. Tan, Q.V. Le, EfficientNetV2: Smaller models and faster training, in: \textit{Proceedings of the 38th International Conference on Machine Learning (ICML)}, PMLR 139, 2021, pp. 10096--10106.

\bibitem{dosovitskiy2021}
A. Dosovitskiy, L. Beyer, A. Kolesnikov, D. Weissenborn, X. Zhai, T. Unterthiner, et al., An image is worth 16x16 words: Transformers for image recognition at scale, in: \textit{International Conference on Learning Representations (ICLR)}, 2021.

\bibitem{liu2021swin}
Z. Liu, Y. Lin, Y. Cao, H. Hu, Y. Wei, Z. Zhang, et al., Swin Transformer: Hierarchical vision transformer using shifted windows, in: \textit{Proceedings of the IEEE/CVF International Conference on Computer Vision (ICCV)}, 2021, pp. 10012--10022.

\bibitem{deng2009}
J. Deng, W. Dong, R. Socher, L.-J. Li, K. Li, L. Fei-Fei, ImageNet: A large-scale hierarchical image database, in: \textit{IEEE Conference on Computer Vision and Pattern Recognition (CVPR)}, 2009, pp. 248--255. \url{https://doi.org/10.1109/CVPR.2009.5206848}.

\bibitem{argenziano1998}
G. Argenziano, G. Fabbrocini, P. Carli, V. De Giorgi, E. Sammarco, M. Delfino, Epiluminescence microscopy for the diagnosis of doubtful melanocytic skin lesions: Comparison of the ABCD rule of dermatoscopy and a new 7-point checklist based on pattern analysis, \textit{Archives of Dermatology} 134 (12) (1998) 1563--1570.

\bibitem{otsu1979}
N. Otsu, A threshold selection method from gray-level histograms, \textit{IEEE Transactions on Systems, Man, and Cybernetics} 9 (1) (1979) 62--66. \url{https://doi.org/10.1109/TSMC.1979.4310076}.

\bibitem{celebi2007}
M.E. Celebi, H.A. Kingravi, B. Uddin, H. Iyatomi, Y.A. Aslandogan, W.V. Stoecker, R.H. Moss, A methodological approach to the classification of dermoscopy images, \textit{Computerized Medical Imaging and Graphics} 31 (6) (2007) 362--373. \url{https://doi.org/10.1016/j.compmedimag.2007.01.003}.

\bibitem{barata2014}
C. Barata, M. Ruela, M. Francisco, T. Mendon{\c{c}}a, J.S. Marques, Two systems for the detection of melanomas in dermoscopy images using texture and color features, \textit{IEEE Systems Journal} 8 (3) (2014) 965--979. \url{https://doi.org/10.1109/JSYST.2013.2271540}.

\bibitem{barata2018}
C. Barata, M.E. Celebi, J.S. Marques, A survey of feature extraction in dermoscopy image analysis of skin cancer, \textit{IEEE Journal of Biomedical and Health Informatics} 23 (3) (2018) 1096--1109. \url{https://doi.org/10.1109/JBHI.2018.2845939}.

\bibitem{tschandl2020}
P. Tschandl, N. Rinner, Z. Apalla, G. Argenziano, N. Codella, A. Halpern, et al., Human--computer collaboration for skin cancer recognition, \textit{Nature Medicine} 26 (8) (2020) 1229--1234. \url{https://doi.org/10.1038/s41591-020-0942-0}.

\bibitem{huang2017}
G. Huang, Z. Liu, L. Van Der Maaten, K.Q. Weinberger, Densely connected convolutional networks, in: \textit{Proceedings of the IEEE Conference on Computer Vision and Pattern Recognition (CVPR)}, 2017, pp. 4700--4708.

\bibitem{sandler2018}
M. Sandler, A. Howard, M. Zhu, A. Zhmoginov, L.-C. Chen, MobileNetV2: Inverted residuals and linear bottlenecks, in: \textit{Proceedings of the IEEE Conference on Computer Vision and Pattern Recognition (CVPR)}, 2018, pp. 4510--4520.

\bibitem{tan2019}
M. Tan, Q. Le, EfficientNet: Rethinking model scaling for convolutional neural networks, in: \textit{International Conference on Machine Learning (ICML)}, PMLR 97, 2019, pp. 6105--6114.

\bibitem{aladhadh2022}
S. Aladhadh, M. Alsanea, M. Aloraini, T. Khan, S. Habib, M. Islam, An effective skin cancer classification mechanism via medical vision transformer, \textit{Sensors} 22 (11) (2022) 4008. \url{https://doi.org/10.3390/s22114008}.

\bibitem{xin2022}
C. Xin, Z. Liu, K. Zhao, L. Miao, Y. Ma, X. Zhu, et al., An improved transformer network for skin cancer classification, \textit{Computers in Biology and Medicine} 149 (2022) 105939. \url{https://doi.org/10.1016/j.compbiomed.2022.105939}.

\bibitem{yap2018}
J. Yap, W. Yolland, P. Tschandl, Multimodal skin lesion classification using deep learning, \textit{Experimental Dermatology} 27 (11) (2018) 1261--1267. \url{https://doi.org/10.1111/exd.13777}.

\bibitem{pacheco2020}
A.G.C. Pacheco, R.A. Krohling, The impact of patient clinical information on automated skin cancer detection, \textit{Computers in Biology and Medicine} 116 (2020) 103545. \url{https://doi.org/10.1016/j.compbiomed.2019.103545}.

\bibitem{ningrum2021}
D.N.A. Ningrum, S.-P. Yuan, W.-M. Kung, C.-C. Wu, I.-S. Tzeng, C.-Y. Huang, et al., Deep learning classifier with patient's metadata of dermoscopic images in malignant melanoma detection, \textit{Journal of Multidisciplinary Healthcare} 14 (2021) 877--885.

\bibitem{bissoto2020}
A. Bissoto, E. Valle, S. Avila, Debiasing skin lesion datasets and models? Not so fast, in: \textit{IEEE/CVF Conference on Computer Vision and Pattern Recognition Workshops (CVPRW)}, 2020, pp. 740--741.

\bibitem{guo2017}
C. Guo, G. Pleiss, Y. Sun, K.Q. Weinberger, On calibration of modern neural networks, in: \textit{Proceedings of the 34th International Conference on Machine Learning (ICML)}, PMLR 70, 2017, pp. 1321--1330.

\bibitem{demvsar2006}
J. Dem{\v{s}}ar, Statistical comparisons of classifiers over multiple data sets, \textit{Journal of Machine Learning Research} 7 (2006) 1--30.

\bibitem{telea2004}
A. Telea, An image inpainting technique based on the fast marching method, \textit{Journal of Graphics Tools} 9 (1) (2004) 23--34. \url{https://doi.org/10.1080/10867651.2004.10487596}.

\bibitem{lin2017}
T.-Y. Lin, P. Goyal, R. Girshick, K. He, P. Doll{\'a}r, Focal loss for dense object detection, in: \textit{Proceedings of the IEEE International Conference on Computer Vision (ICCV)}, 2017, pp. 2980--2988. \url{https://doi.org/10.1109/ICCV.2017.324}.

\bibitem{brier1950}
G.W. Brier, Verification of forecasts expressed in terms of probability, \textit{Monthly Weather Review} 78 (1) (1950) 1--3.

\bibitem{selvaraju2017}
R.R. Selvaraju, M. Cogswell, A. Das, R. Vedantam, D. Parikh, D. Batra, Grad-CAM: Visual explanations from deep networks via gradient-based localization, in: \textit{Proceedings of the IEEE International Conference on Computer Vision (ICCV)}, 2017, pp. 618--626.

\bibitem{combalia2019}
M. Combalia, N.C. Codella, V. Rotemberg, B. Helba, V. Vilaplana, O. Reiter, et al., BCN20000: Dermoscopic lesions in the wild, \textit{arXiv preprint arXiv:1908.02288} (2019).

\bibitem{daneshjou2022}
R. Daneshjou, K. Vodrahalli, R.A. Novoa, M. Jenkins, W. Liang, V. Rotemberg, et al., Disparities in dermatology AI performance on a diverse, curated clinical image set, \textit{Science Advances} 8 (32) (2022) eabq6147. \url{https://doi.org/10.1126/sciadv.abq6147}.

\bibitem{vandermaaten2008}
L. van der Maaten, G. Hinton, Visualizing data using t-SNE, \textit{Journal of Machine Learning Research} 9 (86) (2008) 2579--2605.

\end{thebibliography}

\end{document}
